% ============================================================
\chapter{A2: From Graph to Weights}
\label{chap:sensitivity}
% ============================================================

This chapter presents the allocation machinery (A2, addressing C2) that makes
Q answerable: a family of allocators fed the identical per-window graph, split
into \emph{controls} that can see only the skeleton and \emph{treatments}
that can see the orientations. The family is the construction that
discharges
requirement R2: without it, the value of orientation cannot be measured at
all. Section~\ref{sec:symm} states the problem;
Section~\ref{sec:sem} derives the structural quantities the treatments
consume; Sections~\ref{sec:controls}--\ref{sec:d-family} define the
allocators; Section~\ref{sec:ablation-design} assembles them into the
fixed-graph ablation; and Section~\ref{sec:driver-route} summarises the
driver-based route the project began from, retained as a supporting arm.

\section{The symmetric-clustering problem}
\label{sec:symm}

Hierarchical allocation (Section~\ref{sec:hrp-theme}) consumes a symmetric,
non-negative distance matrix; a causal graph is asymmetric by design, and
this asymmetry is the object of study. The naive fix, symmetrising the
adjacency as
$\tfrac{1}{2}(|M| + |M^\top|)$, discards the very orientations Q asks
about. A more principled alternative also fails: the
bipartite double cover of the directed graph (equivalently the Hermitian
dilation, the symmetric $2N{\times}2N$ block matrix with $M$ and $M^{\top}$
in its off-diagonal blocks) embeds the directed matrix losslessly
in a symmetric object, at the price of doubling the node set, so that every
asset appears once in its source role and once in its sink role.
Hierarchical allocation then produces weights over $2N$ pseudo-assets, and
recombining each asset's two role-weights into a single holding re-poses
precisely the upstream-versus-downstream question the construction was meant
to absorb: the asymmetry is relocated, not resolved. This route was
investigated and abandoned, and it sets the design conclusion for the
chapter: the
symmetric interface cannot be made to carry orientation, so the
allocation steps themselves must change. The challenge is therefore not to
force $M$ into the existing
interface but to build allocation steps that consume the directed
matrix natively, while changing as little else as possible, so that any
performance difference is attributable to orientation and nothing else. The
two structural observations of Section~\ref{sec:hrp-theme} open the way.
First, HRP's dendrogram exists only to produce a leaf \emph{ordering} for
recursive bisection, and a DAG carries a natural ordering for free.
Second, HRP's bisection step consumes cluster \emph{variances} from a
covariance matrix, and a DAG implies one.

\section{The structural model: total effects and the implied covariance}
\label{sec:sem}

Both DYNOTEARS' $W$ and VARLiNGAM's $B_0$, restricted to the asset block and
written in the $i\to j$ convention, define the structural equation model
\begin{equation}
  (I - M^\top)\, x = \varepsilon
  \qquad\Longrightarrow\qquad
  x = (I - M^\top)^{-1} \varepsilon \;=:\; B\,\varepsilon,
  \label{eq:sem}
\end{equation}
where $x$ is the vector of (standardised) asset returns and $\varepsilon$ the
vector of structural shocks. Because the asset block is a DAG, its adjacency
is nilpotent, so the \emph{total-effect matrix}
$B = \sum_{k\ge 0} (M^\top)^k$ terminates in at most $N$ terms and exists
exactly, with no spectral-radius assumption. $B_{ij} = \partial x_i / \partial
\varepsilon_j$ aggregates the influence of a unit shock at asset $j$ on asset
$i$ over all directed paths; it is the Leontief inverse of
input--output economics, and its asymmetry is where orientation enters
allocation natively. (For the non-DAG GRANGER control, $B$ is approximated
by the truncated series with a spectral-radius guard, a documented
approximation.)

The SEM implies a covariance. With $\Sigma_\varepsilon =
\operatorname{diag}(\sigma^2_{\varepsilon,i})$ (diagonality is the SEM's
own independent-shocks assumption, and estimating a dense
$\Sigma_\varepsilon$ would reintroduce sample correlation into a treatment
meant to isolate the graph), the \emph{SEM-implied covariance} in return
units is
\begin{equation}
  \Sigma_{\mathrm{SEM}}
  \;=\; D_\sigma \big( B\, \Sigma_\varepsilon\, B^\top \big) D_\sigma ,
  \label{eq:sigma-sem}
\end{equation}
where the shock variances are estimated from the fit window's structural
residuals $E = X_z (I - M)$, and $D_\sigma$ de-standardises from the
$z$-scored units the graph was fitted in back to return units. The name
merits precision: $\Sigma_{\mathrm{SEM}}$ is the covariance of the
observables that the structural model implies (SEM's $\Sigma(\theta)$,
the reduced form in SVAR terms), whereas the structural shocks' own
covariance is the diagonal $\Sigma_\varepsilon$ above. The
de-standardisation is not optional: recursive bisection compares cluster
variances, and omitting $D_\sigma$ would silently equalise asset
volatilities. Two caveats are stated rather than hidden: the asset-block SEM
marginalises over the drivers, so $\varepsilon$ absorbs driver-explained
variance and its diagonality is an approximation; and
$\Sigma_{\mathrm{SEM}}$ is nearest-PSD projected and lightly ridge-loaded
so dense stress-window graphs cannot produce a singular input.

\section{The controls: correlation and skeleton}
\label{sec:controls}

\paragraph{\HRP{} (the correlation control).} Textbook HRP
\cite{lopezdeprado2016hrp}: the distance $\sqrt{\tfrac{1}{2}(1-\rho_{ij})}$
on the
lookback window's sample correlation, with the sample covariance for
bisection. It is
graph-blind (verified by unit test: changing the graph cannot change its
weights) and anchors the decomposition.

\paragraph{\skelsamp{} (the skeleton control).} The graph replaces the
correlation matrix,
but symmetrically: each asset is embedded by its edge profile
$e_i = [\,M_{i\cdot},\, M_{\cdot i}\,]$ and the clustering distance is
$D_{ij} = \lVert e_i - e_j \rVert_2$, with the sample covariance for
bisection.
This embedding is provably orientation-blind: transposing $M$
(reversing every edge) swaps the two halves of every $e_i$, leaving all
pairwise distances, and hence the dendrogram and the weights, unchanged.
The property is enforced by unit test, so \skelsamp{} is a true skeleton-only
allocator, not merely a weakly-directional one. \skelsamp{} coincides with
the
asset-only Causal-HRP variant of the project's first phase (there called
V0$'$), which allows the new harness to be validated against a
committed historical result (Section~\ref{sec:ablation-design}).

\paragraph{\skelalt{} (a second symmetrisation).} Hierarchical clustering
demands
a symmetric distance and $M$ is not symmetric, so building a skeleton
control
requires choosing a symmetrisation, and the choice is not unique. If
the skeleton rung depended on which one was made, it would be a fact about
the
recipe rather than about the skeleton. \skelalt{} is a control on the
control: the
same pipeline, but with the $\tfrac{1}{2}(|M|+|M^\top|)$ distance, so that
two
assets are close when a strong link joins them rather than when their
edge profiles match. Both are orientation-blind in different
ways, and comparing them tests whether ``the skeleton'' has been
identified accidentally with one particular symmetrisation.

\section{The orientation-aware family}
\label{sec:d-family}

Section~\ref{sec:symm} left hierarchical allocation with exactly two moving
parts, and that fact bounds the design space. The clustering step exists only
to emit an \emph{ordering}; the allocation step consumes only a
\emph{covariance}. Nothing else crosses between them: recursive bisection
receives the ordering and the covariance and no third object, and it splits
the ordered list by position rather than by the dendrogram's own branches. A
causal graph therefore has exactly two doors into HRP, and crossing them
enumerates the family completely (Table~\ref{tab:crossing}).

\begin{table}[h]
  \centering
  \caption{The two doors, crossed. The row is where the ordering comes from,
    the column is what recursive bisection splits on; \skelsamp{} is the control with
    both doors shut to orientation.}
  \label{tab:crossing}
  \begin{tabular}{l l l}
    \toprule
    Ordering from & splits on sample cov. & splits on $\Sigma_{\mathrm{SEM}}$ \\
    \midrule
    clustering on the skeleton & \skelsamp{} (control) & \skelsem{} \\
    topological sort           & \orientsamp{}           & \orientsem{} \\
    \bottomrule
  \end{tabular}
\end{table}

Two points about the rows follow. The first row keeps the clustering step
and feeds
it a graph-derived distance; the second removes the clustering step
entirely, because a DAG already carries an ordering and needs no dendrogram
to produce one. The entry point is therefore ``where the ordering comes
from'', not
``what the clustering step reads''. Each treatment differs from the
control in exactly one respect (\skelsem{} in the covariance alone,
\orientsamp{} in the
ordering alone, \orientsem{} in both), so the results chapter can localise
not just
whether orientation carries value but where. All three are long-only,
fully invested, deterministic and seed-free; all take the identical
$(M,\ \text{returns window})$ inputs as the controls; none has an analogue in
the HRP/HSP literature.

\paragraph{\skelsem{} (SEM-implied covariance).} Clustering is as in
\skelsamp{}: the
ordering is deliberately held at the skeleton, so that
$\skelsem{}-\skelsamp{}$ isolates orientation's effect on the allocation
step
alone. Recursive bisection, however, allocates on $\Sigma_{\mathrm{SEM}}$
of
Eq.~\eqref{eq:sigma-sem} instead of the sample covariance: orientation
enters through the shock-propagation asymmetry of $B$.

\paragraph{\orientsamp{} / \orientsem{} (causal-ordered bisection).}
\label{sec:d2}
The clustering step is removed entirely. HRP's dendrogram exists only to
produce a quasi-diagonalising leaf order, and a DAG supplies an order
natively: assets are sorted upstream$\to$downstream by a topological sort
of
$M$ (ties broken deterministically by total downstream influence
$\sum_r |B_{ri}|$, then name), and the existing recursive bisection
runs on that order, with the sample covariance (\orientsamp{}) or with
$\Sigma_{\mathrm{SEM}}$ (\orientsem{}). This is the most direct response
to the
observation that one cannot cluster on an asymmetric matrix: do not
cluster.
One subtlety, identified by the ablation's own tests, should be recorded:
recursive bisection is mirror-invariant (reversing the leaf order
reproduces the identical partition tree), so \orientsamp{} responds to
orientation
through the partition structure the order induces, not through the
upstream/downstream direction itself.

\section{The fixed-graph ablation}
\label{sec:ablation-design}

The experimental design is what turns the family into an answer to Q.

\paragraph{Design.} At every rebalance the chokepoint
(Section~\ref{sec:asset-graph}) hands the identical graph $M$ to
every allocator in Table~\ref{tab:variants}; everything else (universe,
calendar, costs,
covariance-estimation window, and allocator internals downstream of the
distance/
covariance/ordering choice) is shared code. Orientation is therefore the
sole treatment variable, and the decomposition is two differences with a
common anchor: $\skelsamp{}-\HRP{}$ (skeleton) and
$\orientany{}-\skelsamp{}$ (orientation). The two rungs telescope,
$(\skelsamp{}-\HRP{}) + (\orientany{}-\skelsamp{}) =
\orientany{}-\HRP{}$, so the total causal-graph gain is attributed
without remainder. The reported grid is five graph-consuming allocators $\times$
\{DYNOTEARS, VARLiNGAM\} $\times$ \{189, 252, 378, 504\}-day windows, plus
the magnitude-degraded GRANGER control arm, a sparsity sweep and a cost
sweep. The window axis is a
design variable in its own right: the sample moments every allocator
consumes degrade in opposite directions as the window shrinks (estimation
noise) or grows (staleness across regimes), so tracing the decomposition
over horizons from nine months to two years tests where each kind of
graph information adds value. Every cell of the grid is reported and
priced into the deflation universe
(Section~\ref{sec:believe} discusses the grid's provenance). At the
shortest window the discovery
fit
operates with $n=189$ observations against $d=132$ variables: valid
($n>d$) but thin, a caveat carried through to the results.

\paragraph{Verification before results.} Three gates preceded the grid.
The \emph{unit-test gate}: 60+ tests written first, including the
orientation-blindness of the controls, the orientation-sensitivity of every
treatment, total-effect and SEM-covariance identities against
hand-computed cases, and topological-order properties on random
DAGs. The \emph{cache gate}: all 1,684 (method, window) graph fits
verified
as cache hits with zero refits, so no allocator can have seen a different
graph. The \emph{replication gate}: \skelsamp{} run through the new
harness reproduces
the committed first-phase V0$'$ bundle with $\Delta\text{Sharpe}=0$ and
maximum per-rebalance weight difference $0.0$ across all 215 rebalances,
establishing both that the new machinery is faithful and that
the pipeline is deterministic end-to-end (re-verified on the extended grid
by an exact repeat of a 378-day cell).

\begin{table}[t]
  \centering
  \caption{The allocator family and the supporting driver route. All share
    the universe, calendar, transaction costs and rebalance schedule.}
  \label{tab:variants}
  \begin{tabular}{l l l l}
    \toprule
    Tag & Ordering from & Bisection covariance & Reads orientation? \\
    \midrule
    \HRP{} & correlation distance & sample & no (graph-blind) \\
    \skelsamp{}            & graph embedding      & sample & no (provably) \\
    \skelalt{}           & $\tfrac{1}{2}(|M|{+}|M^\top|)$ & sample & no \\
    \skelsem{}            & graph embedding      & $\Sigma_{\mathrm{SEM}}$ & allocation step \\
    \orientsamp{}            & topological order    & sample & ordering step \\
    \orientsem{}           & topological order    & $\Sigma_{\mathrm{SEM}}$ & both steps \\
    \midrule
    V0/V1/V2      & driver sensitivities & sample & via driver selection \\
    \bottomrule
  \end{tabular}
\end{table}

\section{The incumbent driver route (V0, V1, V2)}
\label{sec:driver-route}

The project began from HSP's driver route, and it is retained as a
supporting
arm, both because it supplies the seed-audit and measurement-problem
evidence of Chapter~\ref{chap:evaluation} and because its failure motivated
Q. Its mechanics are as follows. At each rebalance $K{=}17$ drivers are
selected from the pool, by cumulative correlation with the assets (V0, the
HSP baseline) or by the causal scores of Section~\ref{sec:selection} (V1),
and a per-asset feed-forward network maps the selected drivers to asset
returns; the Jacobian sensitivity vectors $s_i$ embed assets in driver space,
and HRP runs on $\lVert s_i - s_j\rVert_2$ with sample covariance. V2 closes
a feedback loop: the realised holding-period reward, the excess Sharpe over
equal-weight,
\begin{equation}
  R[t] = \mathrm{Sharpe}\big(\text{portfolio over } [t, t{+}21\text{d}]\big)
       - \mathrm{Sharpe}\big(\tfrac{1}{N}\big),
\end{equation}
is attributed to drivers in proportion to $\sum_i |w_i s_{i,d}|$, smoothed by
an EMA $U_d[t] = \gamma\,\mathrm{credit}_d[t] + (1-\gamma)U_d[t-1]$, and
blended into the next selection score
$\alpha\, z(\mathrm{causal}_d) + (1-\alpha)\, z(U_d)$, with lookahead
discipline enforced by keying utility on holding-period-end dates and
a tested leak canary. The update frequency is central: the loop learns from
one
21-observation Sharpe estimate per month, which is the root of the
measurement
problem analysed in Section~\ref{sec:r1}. Unlike the graph-based
allocators, this route is
not deterministic: the network initialisation seed changes the
sensitivities and hence the weights, as quantified by the 20-seed audit of
Section~\ref{sec:seed-audit}.
